% GENERATED FILE. Edit canonical lessons, then run npm run generate:books.
% canonical-source-hash: fnv1a64:26c062b13f60794f
% canonical-lessons: GE-C410-auffahrunfall, GE-C410-auffahren, GE-C410-schleudern, GE-C410-fussgaenger, GE-C410-radfahrer

\chapter{Ein kleiner Unfall}
\label{ch:ge-ein-kleiner-unfall}
\hlchaptermodality{410}

\begin{chapteropening}
\textbf{By the end of this chapter:} \emph{I can describe a small accident: the rear-end collision, driving into a car, skidding, the pedestrian and the cyclist.}

Complete the last lesson of chapter 410: I can describe a small accident: the rear-end collision, driving into a car, skidding, the pedestrian and the cyclist.
\end{chapteropening}

\section[der Auffahrunfall]{der Auffahrunfall — the rear-end collision}

\label{lesson:GE-C410-auffahrunfall}



\begin{quote}
Before the new one: say the German for to close, then the German for the motorway service area.
\end{quote}

\subsection*{You'll want to know: der Auffahrunfall}
\textbf{der Auffahrunfall} — \textquotedblleft{}the rear-end collision\textquotedblright{}.

\begin{quote}
Im Stau gibt es einen kleinen Auffahrunfall.
\end{quote}

\begin{quote}
\emph{There is a small rear-end collision in the traffic jam.}
\end{quote}

Plural: \emph{die Auffahrunfälle}.

\begin{cousinweb}
\emph{Auffahrunfall} is built from \emph{auffahren} 'to drive into the car in front' and \emph{Unfall} 'accident', inherited from Middle High German \emph{unval} 'mishap' (EtymDB).
\end{cousinweb}

\begin{grammarlens}[title={What you've built}]
One more word for telling, explaining and describing at B1.
\end{grammarlens}

\subsection*{Your turn}
\begin{itemize}
\raggedright
  \item \emph{Say it:} \emph{der Auffahrunfall}
  \item \emph{Say it:} \emph{der Auffahrunfall}, once more
  \item \emph{Say it:} a sentence of your own with \emph{der Auffahrunfall}
  \item \emph{Recall:} say the German for to close, then the German for the motorway service area, then say \emph{der Auffahrunfall} again
\end{itemize}

\begin{tcolorbox}[breakable,colback=teal!4,colframe=teal!35!black,title={Before you move on}]
What does der Auffahrunfall mean? (\textquotedblleft{}The rear-end collision\textquotedblright{}.) Say it once more.
\end{tcolorbox}
\section[auffahren]{auffahren — to drive into the back (of another car); to drive up}

\label{lesson:GE-C410-auffahren}



\begin{quote}
Before the new one: say the German for the motorway service area, then the German for the rear-end collision.
\end{quote}

\subsection*{You'll want to know: auffahren}
\textbf{auffahren} — \textquotedblleft{}to drive into the back (of another car); to drive up\textquotedblright{}.

\begin{quote}
Der Fahrer hat nicht gebremst und ist aufgefahren.
\end{quote}

\begin{quote}
\emph{The driver did not brake and drove into the car in front.}
\end{quote}

Separable; Perfekt with \emph{sein}: \emph{er ist auf das Auto aufgefahren}.

\begin{cousinweb}
\emph{auffahren} is built from \emph{auf} 'onto, up' and the inherited verb \emph{fahren} 'to travel, to drive', Proto-Germanic \emph{faraną} (EtymDB).
\end{cousinweb}

\begin{grammarlens}[title={What you've built}]
One more word for telling, explaining and describing at B1.
\end{grammarlens}

\subsection*{Your turn}
\begin{itemize}
\raggedright
  \item \emph{Say it:} \emph{auffahren}
  \item \emph{Say it:} \emph{auffahren}, once more
  \item \emph{Say it:} a sentence of your own with \emph{auffahren}
  \item \emph{Recall:} say the German for the motorway service area, then the German for the rear-end collision, then say \emph{auffahren} again
\end{itemize}

\begin{tcolorbox}[breakable,colback=teal!4,colframe=teal!35!black,title={Before you move on}]
What does auffahren mean? (\textquotedblleft{}To drive into the back (of another car); to drive up\textquotedblright{}.) Say it once more.
\end{tcolorbox}
\section[schleudern]{schleudern — to skid; to fling; to spin (washing)}

\label{lesson:GE-C410-schleudern}



\begin{quote}
Before the new one: say the German for the rear-end collision, then the German for to drive into the back.
\end{quote}

\subsection*{You'll want to know: schleudern}
\textbf{schleudern} — \textquotedblleft{}to skid; to fling; to spin (washing)\textquotedblright{}.

\begin{quote}
Bei Regen schleudert das Auto in der Kurve.
\end{quote}

\begin{quote}
\emph{In the rain the car skids on the bend.}
\end{quote}

Perfekt with \emph{sein} for a car: \emph{das Auto ist geschleudert}. A washing machine also \emph{schleudert}: it spins.

\begin{cousinweb}
\emph{schleudern} is built on \emph{Schleuder} 'sling', first for throwing a stone. EtymDB lists the verb without a source, and the origin of \emph{Schleuder} is not settled.
\end{cousinweb}

\begin{grammarlens}[title={What you've built}]
One more word for telling, explaining and describing at B1.
\end{grammarlens}

\subsection*{Your turn}
\begin{itemize}
\raggedright
  \item \emph{Say it:} \emph{schleudern}
  \item \emph{Say it:} \emph{schleudern}, once more
  \item \emph{Say it:} a sentence of your own with \emph{schleudern}
  \item \emph{Recall:} say the German for the rear-end collision, then the German for to drive into the back, then say \emph{schleudern} again
\end{itemize}

\begin{tcolorbox}[breakable,colback=teal!4,colframe=teal!35!black,title={Before you move on}]
What does schleudern mean? (\textquotedblleft{}To skid; to fling; to spin (washing)\textquotedblright{}.) Say it once more.
\end{tcolorbox}
\section[der Fußgänger]{der Fußgänger — the pedestrian}

\label{lesson:GE-C410-fussgaenger}



\begin{quote}
Before the new one: say the German for to drive into the back, then the German for to skid.
\end{quote}

\subsection*{You'll want to know: der Fußgänger}
\textbf{der Fußgänger} — \textquotedblleft{}the pedestrian\textquotedblright{}.

\begin{quote}
Der Fußgänger geht bei Rot über die Straße.
\end{quote}

\begin{quote}
\emph{The pedestrian crosses the road on a red light.}
\end{quote}

Same form in the plural. A woman: \emph{die Fußgängerin}.

\begin{cousinweb}
\emph{Fußgänger} is inherited from Middle High German \emph{vuozgenger} (EtymDB), built from \emph{Fuß} 'foot' and \emph{Gänger} 'goer', from \emph{gehen}. Dutch \emph{voetganger} is a cousin.
\end{cousinweb}

\begin{grammarlens}[title={What you've built}]
One more word for telling, explaining and describing at B1.
\end{grammarlens}

\subsection*{Your turn}
\begin{itemize}
\raggedright
  \item \emph{Say it:} \emph{der Fußgänger}
  \item \emph{Say it:} \emph{der Fußgänger}, once more
  \item \emph{Say it:} a sentence of your own with \emph{der Fußgänger}
  \item \emph{Recall:} say the German for to drive into the back, then the German for to skid, then say \emph{der Fußgänger} again
\end{itemize}

\begin{tcolorbox}[breakable,colback=teal!4,colframe=teal!35!black,title={Before you move on}]
What does der Fußgänger mean? (\textquotedblleft{}The pedestrian\textquotedblright{}.) Say it once more.
\end{tcolorbox}
\section[der Radfahrer]{der Radfahrer — the cyclist}

\label{lesson:GE-C410-radfahrer}



\begin{quote}
Before the new one: say the German for to skid, then the German for the pedestrian.
\end{quote}

\subsection*{You'll want to know: der Radfahrer}
\textbf{der Radfahrer} — \textquotedblleft{}the cyclist\textquotedblright{}.

\begin{quote}
In der Stadt gibt es viele Radfahrer.
\end{quote}

\begin{quote}
\emph{There are lots of cyclists in the town.}
\end{quote}

Same form in the plural. A woman: \emph{die Radfahrerin}. \emph{Rad fahren} — to cycle.

\begin{cousinweb}
\emph{Radfahrer} is built from \emph{Rad} 'wheel', short for \emph{Fahrrad} 'bicycle', and \emph{Fahrer} 'rider, driver', from the inherited verb \emph{fahren} (EtymDB).
\end{cousinweb}

\begin{grammarlens}[title={What you've built}]
One more word for telling, explaining and describing at B1.
\end{grammarlens}

\subsection*{Your turn}
\begin{itemize}
\raggedright
  \item \emph{Say it:} \emph{der Radfahrer}
  \item \emph{Say it:} \emph{der Radfahrer}, once more
  \item \emph{Say it:} a sentence of your own with \emph{der Radfahrer}
  \item \emph{Recall:} say the German for to skid, then the German for the pedestrian, then say \emph{der Radfahrer} again
\end{itemize}

\begin{tcolorbox}[breakable,colback=teal!4,colframe=teal!35!black,title={Before you move on}]
What does der Radfahrer mean? (\textquotedblleft{}The cyclist\textquotedblright{}.) Say it once more.
\end{tcolorbox}
